\chapter{Planarity and Colouring}

\begin{center}\itshape
A graph can be drawn in the plane without crossings if and only if it contains no subdivision of $K_5$ or $K_{3,3}$.\\
\upshape(Kuratowski, 1930, in our words)
\end{center}

On 23 October 1852 Augustus De Morgan wrote to William Rowan Hamilton: ``A student of mine asked me today to give him a
reason for a fact which I did not know was a fact, and do not yet. He says that if a figure be anyhow divided and the
compartments differently coloured, so that figures with any portion of common boundary line are differently coloured, four
colours may be wanted, but not more.'' The student was Frederick Guthrie, whose brother Francis had noticed it while colouring
a map of the counties of England. Hamilton answered three days later that he was not likely to attempt the ``quaternion of colour''
very soon.

The problem then had one of the most dramatic careers in mathematics. In 1879 Alfred Kempe, a barrister, published a proof;
it was accepted for eleven years, until Percy Heawood found a gap in it, and salvaged from the wreck a proof that \emph{five} colours
suffice. The question stayed open for another 86 years. In 1976 Kenneth Appel and Wolfgang Haken, at the University of
Illinois, settled it with Kempe's own ideas pushed to an extreme: they reduced the problem to $1\,936$ configurations and had a
computer check them, for over a thousand hours. The university's postmark read ``Four colors suffice''. Many mathematicians
were unhappy with a proof that no human could check, and the controversy only calmed down when a simpler computer proof was given
in 1997 and the whole argument was verified in a formal proof assistant by Georges Gonthier in 2005.

This chapter tells the parts of the story that fit in a book. First, what a drawing is, and Euler's formula $n-m+f=2$, which
already shows that $K_5$ and the ``three houses, three utilities'' graph $K_{3,3}$ cannot be drawn without crossings. Then
Kuratowski's theorem, with complete proof: these two graphs are the only obstructions. Then colourings: Heawood's five colour
theorem with Kempe's chains, Brooks' theorem, Vizing's theorem on colouring edges, and the chromatic polynomial, whose value
at $-1$ counts something nobody would expect.

\begin{think}
Three houses must each be connected to water, gas and electricity by pipes that do not cross. Draw as much as you can on
paper. Now try on the surface of a coffee mug, where the handle can carry a pipe over another.
\end{think}

\section{Drawings and Euler's formula}

To prove theorems about drawings, we must say what a drawing is. An \emph{arc} is a polygonal path in $\R^2$, the union of finitely many
segments $[p_0,p_1],[p_1,p_2],\dots,[p_{k-1},p_k]$ meeting only as consecutive segments do; $p_0$ and $p_k$ are its \emph{ends}. A \emph{polygon} is the same with
$p_k=p_0$. A graph is \emph{drawn} in the plane, or is a \emph{plane graph}, when its vertices are distinct points and its edges are arcs joining their
ends, such that the interior of an edge contains no vertex and no point of another edge. A graph is \emph{planar} if it can be drawn.
Restricting to polygonal arcs loses nothing: any drawing with continuous curves can be approximated by a polygonal one, and
F\'ary proved in 1948 that one can even use straight segments. Identify $G$ with the union of its
points and edges. The \emph{faces} of $G$ are the regions of $\R^2\setminus G$: two points are in the same face if an arc avoiding $G$ joins them.
Exactly one face is unbounded, the \emph{outer} face. The \emph{frontier} of a face is the set of points of $G$ that are limits of points of the face.

\begin{lemma}[local structure]\label{lem:local}
Let $G$ be a plane graph.
\begin{enumerate}
\item If $z$ is an interior point of an edge $e$, every small enough disc around $z$ meets $G$ in one or two segments of $e$, and the disc
minus $e$ has two components, the two \emph{sides} of $e$ at $z$. Moving $z$ along $e$, each side stays in one face: so $e$ has a left and a right
side, each contained in a face.
\item If $v$ is a vertex of degree $d\ge1$, every small enough disc around $v$ meets $G$ in $d$ segments from $v$, the beginnings of the edges at $v$, in
distinct directions. They cut the disc into $d$ \emph{sectors} (one if $d=1$). A sector is bordered by two edges consecutive around $v$ and lies in one
face, which contains a side of each of these edges (both sides of the edge if $d=1$).
\item If $G$ has at least one edge, every face contains a side of some edge.
\end{enumerate}
\end{lemma}
\begin{proof}
Everything follows from the finiteness of the drawing. (1) The point $z$ lies in a segment $s$ of $e$, at a positive distance from every
other segment of $G$ not containing $z$; a disc of smaller radius meets only $s$ (or the two segments of $e$ meeting at $z$, if $z$ is a bend),
and a line, or two half-lines from $z$, cut a disc into two parts. Along a segment $s$ of $e$, the points at distance less than $\delta$ on one side of
$s$, away from its ends, form a strip avoiding $G$ for small $\delta$, which joins all the half-discs on that side; at a bend, the two strips
on the same side meet. (2) The same argument at $v$, where the segments leaving $v$ have distinct directions since edges do not overlap; going
around $v$, each sector touches the edges that border it, so it contains a side of each. (3) From a point $x$ of a face, move in a straight
line to a nearest point of $G$; just before arriving, the moving point is in a side of an edge or in a sector at a vertex.
\end{proof}

\begin{lemma}[Jordan curve theorem for polygons]\label{lem:jordan}
A polygon $P$ has exactly two faces, a bounded \emph{interior} and an unbounded \emph{exterior}. Every point of $P$ is in the frontier of both, and the two
sides of each segment of $P$ lie in different faces.
\end{lemma}
\begin{proof}
Regard $P$ as a plane graph, a cycle. By Lemma~\ref{lem:local}, going around $P$, the left side of each edge and the left side of the next one are in one
sector, so all left sides lie in one face; likewise all right sides. By part (3) of that lemma there are at most two faces.

For at least two, fix a direction $u$ not parallel to any segment of $P$. For $x\notin P$, let $\pi(x)$ be the parity of the number of times the half-line
$x+tu$, $t>0$, crosses $P$, where passing through a corner $p_i$ counts as a crossing if the two segments at $p_i$ lie on opposite sides of the line, and
not otherwise. Moving $x$ along a segment avoiding $P$, the half-line sweeps the plane continuously; the number of crossings changes only when the
half-line passes through a corner of $P$, and then by $0$ or $\pm2$ by the rule chosen. So $\pi$ is constant on faces. Two points close to an interior point $z$ of a
segment, on opposite sides of it, have half-lines that differ by exactly one crossing near $z$ (choosing $z$ so that the half-lines meet no corner), so $\pi$
takes both values, and the two sides at $z$ are in different faces. Far away, the half-line misses $P$, so the unbounded face has $\pi=0$ and the face where
$\pi=1$ is bounded. Finally, each point of $P$ is a limit of interior points of segments, near which both faces come arbitrarily close.
\end{proof}

The theorem for arbitrary continuous closed curves, stated by Camille Jordan in 1887, is much harder; it was the first theorem to be
formally verified in a proof assistant at full generality (Hales, 2007). Polygons are all we need.

\begin{lemma}\label{lem:theta}
\begin{enumerate}
\item A plane forest has exactly one face.
\item Let $P_1,P_2,P_3$ be three arcs with the same two ends and otherwise disjoint. Then $P_1\cup P_2\cup P_3$ has exactly three faces, with frontiers
$P_1\cup P_2$, $P_2\cup P_3$ and $P_1\cup P_3$.
\end{enumerate}
\end{lemma}
\begin{proof}
(1) Induction on the number of edges; with no edge, $\R^2$ minus finitely many points is connected. Let $e=uv$ be an edge where $v$ has degree $1$, and $F'$ the
forest without $e$ and $v$, whose complement is connected by induction. By Lemma~\ref{lem:local} the single sector at $v$ contains both sides of $e$, so all
sides of $e$ lie in one face $f_e$ of $F$. The complement of $F'$ is the disjoint union of the open sets $f_e\cup\mathring e\cup\{v\}$ (open, since around points of the
interior $\mathring e$ of $e$ and around $v$ small discs minus $e$ lie in $f_e$) and the union of the other faces of $F$. Being connected, it equals the first one: $f_e$ is the
only face.

(2) Let $u$ be a common end. The three arcs leave $u$ in distinct directions and cut a small disc around $u$ into three sectors, bordered by
$P_1,P_2$, by $P_2,P_3$ and by $P_3,P_1$; each contains one side of each of its two arcs, and a side stays in one face along the whole arc
(Lemma~\ref{lem:local}). Since every face contains a side, every face is the face of one of these three sectors: there are at most three faces.
Two sectors cannot be in the same face: if those bordered by $P_1,P_2$ and by $P_2,P_3$ were, both sides of $P_2$ would be in one face, whereas by
Lemma~\ref{lem:jordan} they are in different faces of the polygon $P_1\cup P_2$, hence of the union. So there are exactly three faces, the one of the sector
bordered by $P_i,P_j$ contains sides of $P_i$ and $P_j$ but not of the third arc, and its frontier is $P_i\cup P_j$.
\end{proof}

\begin{proposition}\label{prop:faces}
In a $2$-connected plane graph every face is bounded by a cycle, and every edge lies in the frontier of exactly two faces.
\end{proposition}
\begin{proof}
By Whitney's theorem (Theorem~\ref{thm:ears}), $G$ has an ear decomposition $H_0\subseteq H_1\subseteq\cdots$. For the cycle $H_0$, this is Lemma~\ref{lem:jordan}.
Suppose it holds for $H$, and add an ear $P$ joining distinct vertices $x,y$ of $H$. The interior of $P$ is connected and avoids $H$, so it lies in a face $f$ of
$H$, bounded by a cycle $C$; the vertices $x$ and $y$ are approached from $f$ along $P$, so they lie on $C$, which they split into two arcs $C_1,C_2$. The face $f$ is a whole
face of the polygon $C$ (a face of $C$ containing $f$ and more would contain a frontier point of $f$, which lies on $C$). By Lemma~\ref{lem:theta}, $C\cup P$ has three faces;
one of them, with frontier $C$, is the face of $C$ not containing $P$; so the other two, with frontiers $C_1\cup P$ and $C_2\cup P$, are the regions into which $P$
cuts $f$. The other faces of $H$ are untouched. So the faces of $H\cup P$ are again bounded by cycles, and the edges of $P$ border exactly two faces, as do those
of $C$ ($f$ was replaced by two faces, one on each side of $C_1$ and of $C_2$).
\end{proof}

\begin{theorem}[Euler's formula]\label{thm:euler2}
A connected plane graph with $n$ vertices, $m$ edges and $f$ faces satisfies $n-m+f=2$.
\end{theorem}
\begin{proof}
Induction on $m-n+1\ge0$. If it is $0$, the graph is a tree and $f=1$ by Lemma~\ref{lem:theta}. Otherwise some edge $e$ lies on a cycle $C$. Its two sides lie in
different faces of $C$ (Lemma~\ref{lem:jordan}), hence in different faces $f_1\ne f_2$ of $G$. In $G-e$, the set $f_1\cup f_2\cup\mathring e$ is open and connected, and avoids
$G-e$; the other faces of $G$ are unchanged. So $G-e$, still connected, has one face fewer and one edge fewer, and $n-m+f$ does not change.
\end{proof}

The formula says that the number of faces does not depend on the drawing. It was discovered by Euler in 1750 for convex polyhedra
(vertices, edges and faces: $8-12+6=2$ for the cube), after Descartes had come close; the polyhedron's surface, punctured at a point of one face
and flattened, is a plane graph. Its consequences are a census of what fits in the plane.

\begin{corollary}\label{cor:planarbounds}
Let $G$ be a simple planar graph with $n\ge3$ vertices and $m$ edges.
\begin{enumerate}
\item $m\le3n-6$, and $m\le2n-4$ if $G$ is $2$-connected and has no triangle.
\item $G$ has a vertex of degree at most $5$.
\item $K_5$ and $K_{3,3}$ are not planar.
\end{enumerate}
\end{corollary}
\begin{proof}
(1) Suppose first $G$ connected, so $m\ge2$. Count the pairs (face, side of an edge in it): there
are $2m$ of them. Every face $\varphi$ contains at least $3$ sides. Indeed $\varphi$ contains a side of some edge $e=uv$; if both $u$ and $v$ had degree $1$, $G$ would be
the single edge $e$. Say $v$ has degree at least $2$: the sector at $v$ containing this side contains a side of the next edge $e'=vw$ around $v$, and $e'\ne e$. If $w$ has
degree $1$, $\varphi$ contains both sides of $e'$; otherwise it contains a side of the next edge $e''$ around $w$, which is not $e$ (else $e$ and $e'$ would both join $v$ and $w$).
So $3f\le2m$, and Euler's formula gives $2=n-m+f\le n-m+\frac23m$, that is $m\le3n-6$. If $G$ is not connected, apply this to each component with at least $3$
vertices; the other components have at most one edge, and adding up gives $m\le3n-6$ again. If $G$ is $2$-connected and has no triangle, each face is bounded by a
cycle of length at least $4$ (Proposition~\ref{prop:faces}) and each edge borders two faces, so $4f\le2m$, and Euler's formula gives $m\le2n-4$.
(2) If every degree were at least $6$, then $2m\ge6n$, against $m\le3n-6$.
(3) $K_5$ has $10>3\cdot5-6$ edges, and $K_{3,3}$, which is $2$-connected without triangles, has $9>2\cdot6-4$ edges.
\end{proof}

\begin{figure}[ht]
\centering
\fig{ch07-k5k33}
\caption{$K_5$ and $K_{3,3}$ drawn with a single crossing; by Corollary~\ref{cor:planarbounds}, no drawing avoids it.}
\label{fig:k5k33}
\end{figure}

\section{Kuratowski's theorem}

Subdividing an edge (replacing it by a path) does not change planarity, so a graph containing a \emph{subdivision} of $K_5$ or $K_{3,3}$ is not planar.
In 1930 Kazimierz Kuratowski proved the converse, which Frink and Smith had also found independently and which Pontryagin may have
known earlier. We follow the proof of Carsten Thomassen (1980), which reduces everything to $3$-connected graphs, where drawings are
almost rigid.

A graph $H$ is a \emph{minor} of $G$ if $G$ has disjoint sets of vertices $V_h$, $h\in V(H)$ (the \emph{branch sets}), each inducing a connected subgraph, with an
edge of $G$ between $V_h$ and $V_{h'}$ whenever $hh'$ is an edge of $H$. Equivalently, $H$ is obtained from a subgraph of $G$ by contracting edges
(and removing loops and parallel edges): contract each branch set to a point. A minor of a minor is a minor, and a subdivision of $H$ in
$G$ gives $H$ as a minor (branch sets: a branch vertex with half of each subdivided edge leaving it).

\begin{lemma}\label{lem:minorsub}
If $G$ has $K_{3,3}$ as a minor, it contains a subdivision of $K_{3,3}$. If $G$ has $K_5$ as a minor, it contains a subdivision of $K_5$ or of $K_{3,3}$.
\end{lemma}
\begin{proof}
For each $h$ choose one edge of $G$ between $V_h$ and $V_{h'}$ for each neighbour $h'$ of $h$, and let $A_h$ be the list of their ends in $V_h$ (with repetitions).
Let $T_h$ be a smallest tree in $G[V_h]$ containing $A_h$. As in the proof of the centroid lemma (Lemma~\ref{lem:centroid}), there is a vertex $c$ of $T_h$ such that each
component of $T_h-c$ contains at most half of the list $A_h$: start anywhere and move towards a component containing more than half; one never moves
back, so this stops.

If $H$ has maximum degree $3$, as $K_{3,3}$ does, each component of $T_h-c$ contains at most one point of $A_h$, so $T_h$ consists of paths from $c$ to the points of
$A_h$ meeting only at $c$. Joining these paths through the chosen edges gives a subdivision of $H$ with branch vertices the points $c$.

For $K_5$ each list has four points. If every $T_h$ is of the same kind, with paths from a centre, we obtain a subdivision of $K_5$. Otherwise some $T_1$ has a
component $K$ of $T_1-c$ containing two points $a_2,a_3$ of the list (towards $V_2$ and $V_3$, say). Let $p\in K$ be the vertex where the paths from $c$ to $a_2$
and to $a_3$ separate. Then $\alpha=$ the vertices of the paths from $p$ to $a_2$ and $a_3$, and $\beta=$ the rest of $T_1$, are connected, disjoint, and adjacent; $\alpha$ is
adjacent to $V_2$ and $V_3$, $\beta$ to $V_4$ and $V_5$. With the branch sets $\{\alpha,V_4,V_5\}$ against $\{\beta,V_2,V_3\}$, $G$ has $K_{3,3}$ as a minor, and the first part
applies.
\end{proof}

So ``no subdivision of $K_5$ or $K_{3,3}$'' and ``no minor $K_5$ or $K_{3,3}$'' are the same condition. We call a graph satisfying it \emph{Kuratowski-free}.

\begin{lemma}[Thomassen, 1980]\label{lem:contractible}
Every $3$-connected graph $G$ with at least $5$ vertices has an edge $xy$ such that $G/xy$ is $3$-connected.
\end{lemma}
\begin{proof}
Suppose not. For each edge $xy$, the graph $G/xy$ has at least $4$ vertices and a separating set of at most $2$ vertices, which must contain the
contracted vertex (otherwise it would separate $G$); so there is a vertex $z$ such that $\{x,y,z\}$ separates $G$. Among all edges $xy$ and such vertices $z$, choose
one for which $G-\{x,y,z\}$ has a component $C$ as small as possible. The vertex $z$ has a neighbour $v$ in $C$, otherwise $\{x,y\}$ would separate $C$ from the rest.
By assumption there is a vertex $w$ such that $\{z,v,w\}$ separates $G$. As $x$ and $y$ are adjacent, they are not in two different components of $G-\{z,v,w\}$, so
some component $D$ of $G-\{z,v,w\}$ contains neither $x$ nor $y$. Again $v$ has a neighbour $u$ in $D$; as $u\notin\{x,y,z\}$ and $v\in C$, $u\in C$. So $D$, which is connected and avoids
$x,y,z$, lies inside $C$, and it does not contain $v$: $|D|<|C|$, against the choice of $C$.
\end{proof}

\begin{lemma}\label{lem:facechoice}
For every face $\varphi$ of a plane graph $G$, $G$ has a drawing whose outer face has the same frontier as $\varphi$.
\end{lemma}
\begin{proof}
Choose a point $p$ in $\varphi$ and apply the inversion $x\mapsto p+\frac{x-p}{|x-p|^2}$, a homeomorphism of $\R^2\setminus\{p\}$ which sends points near $p$ far away and
segments to segments or arcs of circles, preserving angles. The image is a drawing of $G$ by arcs of circles in which $\varphi$ has become the outer face. Replace
each circular arc by an inscribed polygonal arc with very short segments: the edges stay disjoint (distinct edges are at positive distance
away from their common ends, and at a common end they leave in distinct directions, since inversion preserves angles), and the faces keep
their frontiers.
\end{proof}

\begin{lemma}\label{lem:threeconn}
Every $3$-connected Kuratowski-free graph is planar.
\end{lemma}
\begin{proof}
Induction on $n$. For $n=4$, $G=K_4$, which is planar. Let $n\ge5$, take an edge $xy$ as in Lemma~\ref{lem:contractible}, and let $G'=G/xy$, with contracted vertex $v$.
It is $3$-connected and Kuratowski-free (a minor of $G'$ is a minor of $G$), so it has a drawing. The graph $G'-v$ is $2$-connected, so the face $\varphi$ of $G'-v$ containing
the point $v$ is bounded by a cycle $C$ (Proposition~\ref{prop:faces}), and the edges at $v$, which do not meet $G'-v$, run inside $\varphi$ to vertices of $C$. Let
$X=N(x)\setminus\{y\}$ and $Y=N(y)\setminus\{x\}$, both contained in $C$. Let $x_1,\dots,x_k$ be the vertices of $X$ in their order around $C$ ($k\ge2$, as $x$ has degree at least $3$), and
$P_i$ the arc of $C$ from $x_i$ to $x_{i+1}$ (indices modulo $k$). There are three cases (Figure~\ref{fig:thomassen}).

\emph{(a) $Y$ contains three vertices of $X$.} With $x$ and $y$ they are the branch vertices of a subdivision of $K_5$: the three arcs of $C$ between them, and the
edges from $x$ and $y$.

\emph{(b) $Y$ is contained in no $P_i$, and (a) fails.} Then there are $y_1,y_2\in Y$ and $x_i,x_j\in X$ in the order $y_1,x_i,y_2,x_j$ around $C$. (If some $y_1\in Y$ is not in $X$, it lies inside one
$P_i$, and some $y_2\in Y$ lies outside $P_i$; take $x_i,x_{i+1}$. If $Y\subseteq X$, then $Y$ has two vertices, not consecutive among the $x_i$, since $y$ has degree at least $3$
and (a) fails; take one $x_i$ on each side.) The four arcs of $C$ between these vertices, the edges $xx_i$, $xx_j$, $yy_1$, $yy_2$ and $xy$ form a subdivision of $K_{3,3}$
with sides $\{x,y_1,y_2\}$ and $\{y,x_i,x_j\}$.

\emph{(c) $Y\subseteq P_i$ for some $i$.} Both (a) and (b) contradict Kuratowski-freeness, so this case occurs, and we draw $G$. Erase $v$, put $x$ at the
point $v$, and draw the edges $xx_j$ along the old edges $vx_j$. By Lemma~\ref{lem:theta}, applied repeatedly as in Proposition~\ref{prop:faces}, they cut $\varphi$ into faces,
one of which, $\varphi_i$, is bounded by $xx_i\cup P_i\cup x_{i+1}x$. The old edges from $v$ to the vertices of $Y$ inside $P_i$ run in $\varphi_i$. In a small disc around $v$, where all
edges at $v$ are segments from $v$, put $y$ in the sector of $\varphi_i$ next to $v$, join it to $x$ by a segment and to the old edges towards $Y$ by short
segments; for $x_i$ and $x_{i+1}$, if they are in $Y$, run a new edge from $y$ alongside the edge $xx_i$ or $xx_{i+1}$, inside $\varphi_i$. This draws $G$.
\end{proof}

\begin{figure}[ht]
\centering
\fig{ch07-thomassen}
\caption{The three cases of Thomassen's proof. The cycle $C$ bounds the face where $x$ and $y$ must go; the neighbours of $x$ are the dots, those of $y$ the
squares. Left: three common neighbours give $K_5$. Middle: interleaved neighbours give $K_{3,3}$. Right: all neighbours of $y$ in one arc $P_i$, and $y$
can be drawn.}
\label{fig:thomassen}
\end{figure}

\begin{lemma}\label{lem:glue}
Let $G_1$ and $G_2$ be planar graphs such that $V(G_1)\cap V(G_2)$ is a single vertex, or two vertices $x,y$ with $xy$ an edge of both. Then $G_1\cup G_2$ is planar.
\end{lemma}
\begin{proof}
First, contracting an edge $ab$ of a plane graph gives a plane graph: erase $b$, and redraw each other edge $bc$ as an arc that starts at $a$, runs alongside $ab$
on the side where $bc$ leaves $b$ (at slightly different distances for different edges), turns around $b$ inside a small disc and continues along $bc$.
Now, by Lemma~\ref{lem:facechoice}, draw $G_1$ with the common vertex $s$, or the common edge $xy$, on the frontier of the outer face, and draw $G_2$ likewise, far away,
with its copy $s'$, or $x'y'$, of the common part. In the first case join $s$ to $s'$ by an arc through the outer faces. In the second, take an arc $Q$ in the
outer faces from a point beside the middle of $xy$ to a point beside the middle of $x'y'$, and join $x$ to $x'$ by an arc running alongside $xy$, then alongside
$Q$ on its left, then alongside $x'y'$; and $y$ to $y'$ in the same way on the right of $Q$. Contracting the new edges (and merging the two copies of $xy$)
gives a drawing of $G_1\cup G_2$.
\end{proof}

\begin{theorem}[Kuratowski, 1930; Wagner, 1937]\label{thm:kuratowski}
For a graph $G$, the following are equivalent: (i) $G$ is planar; (ii) $G$ contains no subdivision of $K_5$ or $K_{3,3}$; (iii) $G$ has no minor $K_5$ or $K_{3,3}$.
\end{theorem}
\begin{proof}
(i)$\Rightarrow$(ii): a drawing of a subdivision of $K_5$ or $K_{3,3}$ would give one of $K_5$ or $K_{3,3}$, merging each subdivided edge into one arc, against
Corollary~\ref{cor:planarbounds}. (ii)$\Rightarrow$(iii) is Lemma~\ref{lem:minorsub}. For (iii)$\Rightarrow$(i) we induct on $n$; graphs with at most $4$ vertices are subgraphs of $K_4$, which is planar. Let $G$ be
Kuratowski-free with $n\ge5$; its subgraphs are Kuratowski-free too. If $G$ is disconnected, draw its components by induction, far apart. If $G$ has a cut
vertex $s$, it is the union of two smaller graphs meeting only at $s$, planar by induction, and Lemma~\ref{lem:glue} applies. If $G$ is $3$-connected,
Lemma~\ref{lem:threeconn} applies. Otherwise $G$ is $2$-connected with a separating pair $\{x,y\}$: $G=G_1\cup G_2$ with $V(G_1)\cap V(G_2)=\{x,y\}$, both $G_i-\{x,y\}$ non-empty,
and no edge between them. Let $G_i^+=G_i+xy$ (adding the edge if it is missing). Each $G_i^+$ is Kuratowski-free. Indeed, a subdivision $K$ of $K_5$ or $K_{3,3}$ in $G_1^+$
would use the edge $xy$, as $G_1\subseteq G$; since $G$ is $2$-connected, $x$ and $y$ both have neighbours in each component of $G-\{x,y\}$ (otherwise the other one would be a cut
vertex), so there is a path from $x$ to $y$ through $G_2-\{x,y\}$, which meets $K$ only at $x$ and $y$; replacing the edge $xy$ by this path gives a subdivision of $K_5$ or $K_{3,3}$
in $G$, which is impossible. The graphs $G_i^+$ have fewer vertices than $G$, so they are planar by induction, and by Lemma~\ref{lem:glue} so is $G_1^+\cup G_2^+$, which
contains $G$.
\end{proof}

Hopcroft and Tarjan found in 1974 an algorithm that decides planarity in linear time, and finds a drawing; Tome~C2 explains the
simpler left-right algorithm of de Fraysseix and Rosenstiehl, which the notebook uses as an oracle. Wagner's form of the theorem has a
majestic generalisation: Robertson and Seymour proved, in twenty-three papers written from 1983 to 2004, that every family of graphs
closed under taking minors is characterised by a \emph{finite} list of forbidden minors. For graphs drawable on the torus, thousands of forbidden minors
are known, and nobody is sure that the list is complete.

\section{Colouring vertices}

A \emph{proper colouring} of $G$ with $k$ colours gives each vertex one of $k$ colours so that adjacent vertices differ; the \emph{chromatic number} $\chi(G)$ is the
least such $k$. Colourings model conflicts: exams sharing a student cannot be at the same time, two variables alive at the same moment of
a program cannot share a register, two radio transmitters close to each other need different frequencies. Deciding whether $\chi(G)\le3$ is
NP-complete (Tome~C3), so we look for bounds. The \emph{clique number} $\omega(G)$, the size of a largest complete subgraph, is a lower bound; it can be far
from the truth, since there are triangle-free graphs of any chromatic number (Exercise~\ref{ex:mycielski}).

\begin{proposition}[greedy colouring]\label{prop:greedy}
Colour the vertices in some order $v_1,\dots,v_n$, giving each the smallest colour not used by its already coloured neighbours. If each $v_i$ has at most $d$
neighbours among $v_1,\dots,v_{i-1}$, at most $d+1$ colours are used. In particular $\chi(G)\le\Delta(G)+1$, and $\chi(G)\le d+1$ if every subgraph of $G$ has a vertex of
degree at most $d$ ($G$ is \emph{$d$-degenerate}).
\end{proposition}
\begin{proof}
When $v_i$ is coloured, at most $d$ colours are forbidden, so one of $1,\dots,d+1$ is free. For a $d$-degenerate graph, choose $v_n$ of degree at most $d$, then $v_{n-1}$ of
degree at most $d$ in $G-v_n$, and so on: each $v_i$ has at most $d$ neighbours before it.
\end{proof}

By Corollary~\ref{cor:planarbounds}, planar graphs are $5$-degenerate (subgraphs of planar graphs are planar), so they are $6$-colourable. Kempe's idea
improves this to five. A \emph{Kempe chain} for colours $\alpha,\beta$ is a connected component of the subgraph spanned by the vertices of colours $\alpha$ and $\beta$; exchanging
$\alpha$ and $\beta$ on a Kempe chain gives another proper colouring.

\begin{theorem}[Heawood, 1890]\label{thm:fivecolour}
Every planar graph is $5$-colourable.
\end{theorem}
\begin{proof}
Induction on $n$; we may assume $G$ simple and drawn. Take a vertex $v$ of degree at most $5$, and colour $G-v$ with five colours. If the neighbours of $v$
use at most four colours, $v$ gets the fifth. Otherwise $v$ has five neighbours $v_1,\dots,v_5$, in this order around $v$, with colours $1,\dots,5$. Consider the
Kempe chain $K$ of colours $1,3$ containing $v_1$. If it does not contain $v_3$, exchange $1$ and $3$ on $K$: now $v_1$ has colour $3$, colour $1$ is free around $v$, and
$v$ gets it. Otherwise there is a path from $v_1$ to $v_3$ coloured $1$ and $3$, which with $v$ forms a cycle $C$ (Figure~\ref{fig:kempe}). Near $v$, the edges $vv_2$ and $vv_4$
leave $v$ in the two different sectors cut by $vv_1$ and $vv_3$, which contain the two sides of the edge $vv_1$; by Lemma~\ref{lem:jordan}, $v_2$ and $v_4$ lie in different
faces of $C$ (neither is on $C$, whose vertices are $v$ and vertices coloured $1$ or $3$). Any path from $v_2$ to $v_4$ meets $C$, at a vertex, which is not coloured $2$
or $4$. So the Kempe chain of colours $2,4$ containing $v_2$ does not contain $v_4$; exchanging $2$ and $4$ on it frees colour $2$ for $v$.
\end{proof}

\begin{figure}[ht]
\centering
\fig{ch07-kempe}
\caption{Kempe chains around a vertex $v$ of degree $5$. If a chain of colours $1$ and $3$ joins $v_1$ to $v_3$, it closes a cycle with $v$ that separates $v_2$ from $v_4$,
so the chain of colours $2$ and $4$ starting at $v_2$ cannot reach $v_4$.}
\label{fig:kempe}
\end{figure}

Kempe's 1879 argument for four colours ran the same way at a vertex of degree $5$, with two exchanges at once; Heawood showed that the two
exchanges can interfere. The four colour theorem is true, but its proof needs, instead of one configuration (a vertex of degree at most
$5$), a list of $633$ ``reducible'' configurations, and a discharging argument showing that every planar triangulation contains one; both
halves are checked by computer (Robertson, Sanders, Seymour and Thomas, 1997).

For general graphs, the greedy bound $\Delta+1$ is attained by complete graphs and odd cycles. These are the only exceptions.

\begin{theorem}[Brooks, 1941]\label{thm:brooks}
If $G$ is connected and is neither a complete graph nor an odd cycle, then $\chi(G)\le\Delta(G)$.
\end{theorem}
\begin{proof}[Proof (Lov\'asz, 1975)]
Let $\Delta=\Delta(G)$. If $\Delta\le2$, $G$ is a path or an even cycle, $2$-colourable. Let $\Delta\ge3$. The trick is to colour greedily in an order in which every vertex but the
last has a neighbour coming later (so at most $\Delta-1$ neighbours earlier), and to treat the last vertex specially.

\emph{$G$ is not regular.} Let $r$ be a vertex of degree less than $\Delta$, and order the vertices by decreasing distance from $r$, ending with $r$. Each vertex other than $r$
has a later neighbour on a shortest path to $r$, and $r$ has fewer than $\Delta$ neighbours: $\Delta$ colours suffice.

\emph{$G$ is regular with a cut vertex $s$.} For each component $D$ of $G-s$, the graph $G[D\cup\{s\}]$ is connected and not regular ($s$ has smaller degree there), so it is
$\Delta$-colourable; permuting colours, all give $s$ the same colour, and they combine.

\emph{$G$ is regular and $2$-connected.} Suppose we find a vertex $z$ with two non-adjacent neighbours $x,y$ such that $G-\{x,y\}$ is connected. Colour $x$ and $y$ with colour $1$;
then colour the vertices of $G-\{x,y\}$ by decreasing distance from $z$ in $G-\{x,y\}$, ending with $z$. Every vertex but $z$ has a later neighbour, and $z$ has two
neighbours of the same colour, so at most $\Delta-1$ colours around it: $\Delta$ colours suffice. To find $z,x,y$: if $G$ is $3$-connected, take any vertex $z$ with two non-adjacent
neighbours (one exists, since $G$ is connected and not complete). Otherwise some $\{z,w\}$ separates $G$, and $H=G-z$ is connected with a cut vertex. Choose a cut vertex $w_1$ of
$H$ and a component $D_1$ of $H-w_1$ with $|D_1|$ smallest; then no vertex $u$ of $D_1$ is a cut vertex of $H$ (a component of $H-u$ avoiding $w_1$ would lie in $D_1\setminus\{u\}$). Next
choose a cut vertex $w_2$ and a component $D_2$ of $H-w_2$ disjoint from $D_1\cup\{w_1\}$ with $|D_2|$ smallest (the other components of $H-w_1$ qualify); by the same argument no
vertex of $D_2$ is a cut vertex of $H$ (a component of $H-u$ avoiding $w_2$ lies in $D_2$, hence avoids $D_1\cup\{w_1\}$). The vertex $z$ has a neighbour $x\in D_1$, otherwise $w_1$ would
separate $D_1$ in $G$; likewise a neighbour $y\in D_2$. They are not adjacent, as $D_1$ is a component of $H-w_1$ and $y\notin D_1\cup\{w_1\}$. Finally $H-\{x,y\}$ is connected: every vertex of
$D_1-x$ reaches $w_1$ in $H-x$ without leaving $D_1\cup\{w_1\}$, and every other vertex reaches $w_1$ in $H-y$ before entering $D_1$. As $z$ has $\Delta\ge3$ neighbours, it has one
besides $x$ and $y$, and $G-\{x,y\}$ is connected.
\end{proof}

\section{Colouring edges}

An \emph{edge colouring} gives colours to edges so that edges sharing a vertex differ; the least number of colours is the \emph{chromatic index} $\chi'(G)$. Each
colour class is a matching, and the $\Delta$ edges at a vertex of maximum degree need distinct colours, so $\chi'(G)\ge\Delta$. For bipartite graphs, K\H onig's
theorem $\chi'=\Delta$ (Exercise~\ref{ex:edgecolour}) settles the matter; in general one more colour always suffices.

\begin{theorem}[Vizing, 1964]\label{thm:vizing}
Every simple graph has $\chi'(G)\le\Delta(G)+1$.
\end{theorem}
\begin{proof}
Induction on the number of edges; let $e=xy_0$ be an edge, and colour $G-e$ with $\Delta+1$ colours. Every vertex has degree at most $\Delta$, so it \emph{misses} some colour (no
edge of that colour at it). Let $\alpha$ be a colour missing at $x$. A \emph{fan} is a sequence $y_0,y_1,\dots,y_k$ of distinct neighbours of $x$ such that the colour of $xy_i$ is missing at
$y_{i-1}$, for $1\le i\le k$ (Figure~\ref{fig:vizing}); take a longest one. For $0\le j\le k$, \emph{shifting the fan at $j$} means giving $xy_i$ the colour of $xy_{i+1}$ for $i<j$ and
uncolouring $xy_j$: this is again a proper colouring of $G$ minus one edge, since the colour of $xy_{i+1}$ is missing at $y_i$, and at $x$ the colours are only moved around.

Let $\beta$ be a colour missing at $y_k$. If $\beta$ is also missing at $x$, shift at $k$ and colour $xy_k$ with $\beta$: done. Otherwise some edge $xy$ has colour $\beta$; since the fan
cannot be extended, $y$ is one of the $y_i$, with $1\le i<k$ ($xy_0$ is uncoloured, and $\beta$ is missing at $y_k$).

Shift at $k$. Now $xy_{i-1}$ has colour $\beta$, $y_i$ misses $\beta$ (its $\beta$-edge $xy_i$ was recoloured), and $x$ still misses $\alpha$. If $y_k$ also misses $\alpha$, colour $xy_k$ with $\alpha$:
done. Otherwise let $P$ be the maximal path starting at $y_k$ whose edges are coloured alternately $\alpha$ and $\beta$, beginning with $\alpha$ (the edges of two colours form paths and
cycles, and $y_k$ misses $\beta$). If $P$ does not reach $x$, exchange $\alpha$ and $\beta$ on $P$: then $y_k$ misses $\alpha$, $x$ still misses it, and $xy_k$ gets $\alpha$. If $P$ reaches $x$, it does
so through its only $\alpha$- or $\beta$-edge, $xy_{i-1}$, and ends there, since $x$ misses $\alpha$.

In this last case, go back to the colouring shifted at $i-1$ instead. It differs from the colouring shifted at $k$ only on edges at $x$: in it, $xy_{i-1}$ is uncoloured
and $xy_i$ has colour $\beta$. The path $P$ without its last edge is an $\alpha\beta$-path from $y_k$ to $y_{i-1}$; it is still a maximal one, as $y_{i-1}$ misses $\beta$ (by the fan property)
and the only new $\beta$-edge, $xy_i$, has its ends off this path ($y_i$ could only have been an end of $P$, having missed $\beta$). Exchange $\alpha$ and $\beta$ on it: $y_{i-1}$ now misses
$\alpha$, and so does $x$; colour $xy_{i-1}$ with $\alpha$.
\end{proof}

\begin{figure}[ht]
\centering
\fig{ch07-vizing}
\caption{A fan at $x$ in Vizing's proof: the colour of each edge $xy_i$ is missing at the previous vertex $y_{i-1}$, so the colours can be shifted down the fan,
freeing $xy_j$ for any $j$.}
\label{fig:vizing}
\end{figure}

So every graph is of \emph{class 1} ($\chi'=\Delta$) or \emph{class 2} ($\chi'=\Delta+1$). Deciding which is NP-complete, even for $\Delta=3$ (Holyer, 1981), although
almost all graphs are of class 1. A regular graph with an odd number of vertices is of class 2 (Exercise~\ref{ex:classtwo}), and so is the
Petersen graph, which explains why the four colour theorem is equivalent to the $3$-edge-colourability of bridgeless cubic planar graphs (Tait, 1880):
the Petersen graph is not planar.

\section{The chromatic polynomial}

Let $P(G,k)$ be the number of proper colourings of $G$ with colours $1,\dots,k$. For the empty graph on $n$ vertices it is $k^n$, for $K_n$ it is $k(k-1)\cdots(k-n+1)$, for a tree
$k(k-1)^{n-1}$ (colour a root, then each vertex differently from its parent). George Birkhoff introduced it in 1912, hoping to prove the four
colour theorem by showing $P(G,4)>0$ for planar $G$ with analysis; this never worked, but the polynomial is beautiful. For an edge $e=uv$, write $G-e$ for the
deletion and $G/e$ for the contraction, with parallel edges merged.

\begin{proposition}[deletion--contraction]\label{prop:delcon}
$P(G,k)=P(G-e,k)-P(G/e,k)$. Consequently $P(G,k)$ is a polynomial in $k$ of degree $n$ with leading coefficient $1$, its next coefficient is $-m$, and its
coefficients alternate in sign.
\end{proposition}
\begin{proof}
The colourings of $G-e$ split into those with $u,v$ coloured differently, which are the colourings of $G$, and those with the same colour, which are the
colourings of $G/e$. Induction on $m$, with the empty graph as base, gives the rest: $G-e$ has $n$ vertices and $m-1$ edges, $G/e$ has $n-1$ vertices and at most $m-1$
edges, and $(-1)^{n-i}$ times the coefficient of $k^i$ is non-negative in both, so subtracting keeps the alternation; the coefficient of $k^{n-1}$ is $-(m-1)-1$.
\end{proof}

The coefficients themselves count something. Number the edges $e_1<e_2<\dots<e_m$. A \emph{broken circuit} is a cycle with its smallest edge removed.

\begin{theorem}[Whitney, 1932]\label{thm:whitney}
$P(G,k)=\sum_{i\ge0}(-1)^ia_ik^{n-i}$, where $a_i$ is the number of sets of $i$ edges containing no broken circuit.
\end{theorem}
\begin{proof}
For a set $A$ of edges, the colourings in which every edge of $A$ is monochromatic are constant on the components of the graph $(V,A)$; there are $k^{c(A)}$ of them,
$c(A)$ being the number of components. By inclusion--exclusion over the set of monochromatic edges,
\[
P(G,k)=\sum_{A\subseteq E}(-1)^{|A|}k^{c(A)}.
\]
Call an edge $e$ \emph{active} for $A$ if $e$ is the smallest edge of a cycle $C$ with $C-e\subseteq A$; then $A$ contains a broken circuit if and only if it has an active edge. For such $A$, let
$e$ be its smallest active edge and $\iota(A)=A\triangle\{e\}$. Adding or removing $e$ does not change $c(A)$, since the ends of $e$ are joined by $C-e\subseteq A\setminus\{e\}$, and it
changes the sign. And $e$ is still the smallest active edge for $\iota(A)$: it is active ($C-e$ is untouched); and if $f<e$ were active for $\iota(A)$, through a cycle $C'$, then $C'-f\not\subseteq A$
(otherwise $f$ would be active for $A$), so $e\in C'-f$, and the ends of $f$ are joined by $(C'-f-e)\cup(C-e)\subseteq A$, giving a cycle $C''$ through $f$ in $A\cup\{f\}$ whose edges, from $C'$ or $C$, are all
at least $f$; so $f$ would be active for $A$. Thus $\iota$ is an involution pairing the sets with a broken circuit, with opposite signs and equal $k^{c(A)}$: they cancel. The
remaining sets contain no cycle (a cycle contains a broken circuit), so they are forests with $c(A)=n-|A|$.
\end{proof}

\begin{theorem}[Stanley, 1973]\label{thm:stanley}
$(-1)^nP(G,-1)$ is the number of acyclic orientations of $G$.
\end{theorem}
\begin{proof}
Let $a(G)$ be the number of acyclic orientations. Both sides equal $1$ for graphs without edges, and $(-1)^nP(G,-1)$ satisfies $(-1)^nP(G,-1)=(-1)^nP(G-e,-1)+(-1)^{n-1}P(G/e,-1)$
by Proposition~\ref{prop:delcon}. So it suffices to show $a(G)=a(G-e)+a(G/e)$ for $e=uv$, and induct on $m$.

Take an acyclic orientation $O$ of $G-e$. Orienting $e$ as $u\to v$ creates a directed cycle only if $O$ has a directed path from $v$ to $u$, and $v\to u$ only if it has one from $u$
to $v$; both cannot happen, as $O$ is acyclic. So $O$ extends to one or two acyclic orientations of $G$, two exactly when $O$ has no directed path between $u$ and $v$. Hence
$a(G)=a(G-e)+b$, where $b$ is the number of such $O$. These $O$ correspond to the acyclic orientations of $G/e$: if $w$ is adjacent to both $u$ and $v$, the edges $uw$ and $vw$ point
the same way in $O$ (otherwise there would be a path $u\to w\to v$ or $v\to w\to u$), so $O$ induces an orientation of $G/e$, acyclic since a directed cycle through the merged vertex
would lift to a directed path between $u$ and $v$ or a cycle in $O$; conversely an acyclic orientation of $G/e$ lifts to an acyclic $O$ without $u$--$v$ paths, which would
project to cycles. So $b=a(G/e)$.
\end{proof}

For $K_n$, $(-1)^nP(K_n,-1)=n!$: the acyclic orientations of a complete graph are the linear orders, the tournaments without cycles of
Chapter~5. ``Evaluating a counting polynomial at a negative number counts something else'' is a phenomenon called \emph{combinatorial reciprocity},
and it has become a theory of its own (Beck and Sanyal's book).

\section{Beyond the plane}

On a surface with $g$ handles (the sphere has $g=0$, the torus and the coffee mug $g=1$), Euler's formula becomes $n-m+f=2-2g$ for suitable drawings,
so $m\le3n-6+6g$. Applied to a smallest graph that needs more than $h$ colours, whose minimum degree is at least $h$ and which has more than $h$ vertices,
this counting gives Heawood's bound of 1890,
\[
\chi\le\Bigl\lfloor\frac{7+\sqrt{1+48g}}2\Bigr\rfloor\qquad(g\ge1).
\]
For the torus it gives $7$, and $K_7$ can indeed be drawn on a torus: seven countries on a doughnut can all touch each other. Ringel
and Youngs proved in 1968 that the bound is attained on every surface with $g\ge1$; for the sphere ($g=0$) it would say $4$, and that is the four colour
theorem, the only case the surface methods could not reach. On the coffee mug of the opening question the three houses can be connected:
$K_{3,3}$, and even $K_7$, can be drawn on a torus without crossings.

\begin{beauty}
Euler's formula is a statement about a drawing that does not depend on the drawing: the alternating count $n-m+f$ is an invariant of the
surface. It is the first instance of the Euler characteristic, which topology later computed for every space through homology, and which
reappears in the Gauss--Bonnet theorem as the total curvature of a surface divided by $2\pi$. A child's count of corners, edges and faces of a box
is the shadow of one of the deepest invariants of mathematics.
\end{beauty}

\section{Exercises}

\begin{exercise}[$\star$]\label{ex:girth}
Let $G$ be a $2$-connected plane graph with $n$ vertices, $m$ edges, and no cycle of length less than $g\ge3$. Show that $m\le\frac g{g-2}(n-2)$. Deduce that the
Petersen graph is not planar.
\end{exercise}

\begin{exercise}[$\star\star$]\label{ex:platonic}
\emph{The Platonic solids.} A convex polyhedron is regular if all vertices have the same degree $k\ge3$ and all faces have the same number $d\ge3$ of sides. Using
Euler's formula, show that $(k,d)$ is one of $(3,3)$, $(3,4)$, $(4,3)$, $(3,5)$, $(5,3)$, and find $n$, $m$, $f$ in each case.
\end{exercise}

\begin{exercise}[$\star$]\label{ex:football}
A football (or a fullerene molecule) is a convex polyhedron whose faces are pentagons and hexagons, with three faces at each vertex. Show that it has
exactly $12$ pentagons, however many hexagons it has.
\end{exercise}

\begin{exercise}[$\star\star$]\label{ex:crossing}
The \emph{crossing number} $\mathrm{cr}(G)$ is the least number of crossings in a drawing of $G$ where edges may cross (at single points, two at a time). Show that
$\mathrm{cr}(G)\ge m-3n+6$ for $n\ge3$, and that $\mathrm{cr}(K_6)=3$.
\end{exercise}

\begin{exercise}[$\star\star$]\label{ex:hakimi}
Using Hall's theorem (Theorem~\ref{thm:hall}), show that the edges of every simple planar graph can be oriented so that every vertex has out-degree at most $3$.
\end{exercise}

\begin{exercise}[$\star$]\label{ex:greedyorder}
Show that for every graph there is an order in which greedy colouring uses exactly $\chi(G)$ colours, and that for the graph $K_{n,n}$ minus a perfect matching
there is an order in which it uses $n$ colours.
\end{exercise}

\begin{exercise}[$\star$]\label{ex:interval}
\emph{Register allocation.} An \emph{interval graph} has a vertex for each of finitely many intervals of $\R$, adjacent when the intervals intersect (variables alive
during intervals of a program, adjacent when alive together). Show that greedy colouring in the order of left endpoints uses $\omega(G)$ colours.
\end{exercise}

\begin{exercise}[$\star\star$]\label{ex:mycielski}
\emph{Mycielski, 1955.} From a graph $G$ with vertices $v_1,\dots,v_n$ build $M(G)$ by adding vertices $u_1,\dots,u_n,w$, with $u_i$ adjacent to the neighbours of $v_i$ and $w$
adjacent to all $u_i$. Show that if $G$ has no triangle, neither has $M(G)$, and that $\chi(M(G))=\chi(G)+1$. Deduce that there are triangle-free graphs of every chromatic
number.
\end{exercise}

\begin{exercise}[$\star$]\label{ex:classtwo}
Show that a $\Delta$-regular graph with an odd number of vertices has chromatic index $\Delta+1$. What does this say about $K_n$ for $n$ odd?
\end{exercise}

\begin{exercise}[$\star$]\label{ex:chromcycle}
Show that $P(C_n,k)=(k-1)^n+(-1)^n(k-1)$ for the cycle $C_n$, $n\ge3$. Check Stanley's theorem on it.
\end{exercise}

\begin{exercise}[$\star$]\label{ex:eulerc}
Show that a plane graph with $n$ vertices, $m$ edges, $f$ faces and $c$ components satisfies $n-m+f=1+c$.
\end{exercise}

\begin{exercise}[$\star\star\star$]\label{ex:listfive}
\emph{Thomassen's five list colouring, 1994.} Let $G$ be a plane graph whose inner faces are triangles and whose outer face is bounded by a cycle $C=v_1v_2\cdots v_k$. Suppose
$v_1,v_2$ are coloured differently, every other vertex of $C$ has a list of at least $3$ allowed colours, and every inner vertex a list of at least $5$. Show that the colouring
extends, every vertex getting a colour from its list. (Hence planar graphs can be coloured from any lists of size $5$, a strengthening of Heawood's theorem.)
\end{exercise}

\section{Hints}

\begin{hint}{ex:girth}
Each face has at least $g$ edges on its boundary cycle, and each edge borders two faces.
\end{hint}
\begin{hint}{ex:platonic}
$kn=2m=df$; divide Euler's formula by $2m$.
\end{hint}
\begin{hint}{ex:football}
Count edges twice: $3n=2m=5p+6h$.
\end{hint}
\begin{hint}{ex:crossing}
Remove one edge per crossing. For $K_6$, nest two triangles.
\end{hint}
\begin{hint}{ex:hakimi}
Match each edge with one of three copies of one of its ends.
\end{hint}
\begin{hint}{ex:greedyorder}
List the colour classes of an optimal colouring one after the other.
\end{hint}
\begin{hint}{ex:interval}
When an interval is coloured, its earlier neighbours all contain its left endpoint.
\end{hint}
\begin{hint}{ex:mycielski}
Given a colouring of $M(G)$, recolour each $v_i$ that has $w$'s colour with the colour of $u_i$.
\end{hint}
\begin{hint}{ex:classtwo}
How many edges can one colour class have?
\end{hint}
\begin{hint}{ex:chromcycle}
Delete and contract an edge of $C_n$.
\end{hint}
\begin{hint}{ex:eulerc}
Follow the proof of Theorem~\ref{thm:euler2}.
\end{hint}
\begin{hint}{ex:listfive}
Induction on $n$. Use a chord of $C$ if there is one; otherwise remove $v_k$ and protect it with two colours.
\end{hint}

\section{Solutions}

\begin{solution}{ex:girth}
\idea Double counting of edge--face incidences, as for $K_{3,3}$, with $g$ in place of $4$.
\writeup By Proposition~\ref{prop:faces}, each face is bounded by a cycle, of length at least $g$, and each edge borders exactly two faces: $gf\le2m$. Euler's formula gives
$2=n-m+f\le n-m+\frac{2m}g$, so $m(1-\frac2g)\le n-2$. The Petersen graph is $3$-connected, has $n=10$, $m=15$ and girth $5$; the bound would give $15\le\frac53\cdot8$, false.
\end{solution}

\begin{solution}{ex:platonic}
\idea Euler's formula becomes an inequality between $\frac1k$ and $\frac1d$ with few solutions.
\writeup From $kn=2m=df$ and $n-m+f=2$: $\frac{2m}k-m+\frac{2m}d=2$, so $\frac1k+\frac1d=\frac12+\frac1m>\frac12$. With $k,d\ge3$, the solutions are $(3,3)$, $(3,4)$, $(4,3)$, $(3,5)$, $(5,3)$ (if
both are at least $4$, or one is $3$ and the other at least $6$, the sum is at most $\frac12$). Then $m$ follows from $\frac1m=\frac1k+\frac1d-\frac12$: $6$, $12$, $12$, $30$, $30$. Tetrahedron
$(4,6,4)$, cube $(8,12,6)$, octahedron $(6,12,8)$, dodecahedron $(20,30,12)$, icosahedron $(12,30,20)$, listing $(n,m,f)$.
\end{solution}

\begin{solution}{ex:football}
\idea Written in terms of faces, Euler's formula gives each hexagon weight $0$ and each pentagon weight $1$.
\writeup Let $p$ and $h$ be the numbers of pentagons and hexagons. Each edge borders two faces and each vertex three: $2m=5p+6h$ and $3n=2m$. Euler's formula times $6$:
$6n-6m+6f=12$, with $6n=4m$: $-2m+6p+6h=12$, so $-(5p+6h)+6p+6h=12$ and $p=12$. The classic football has $12$ pentagons and $20$ hexagons, the molecule $C_{60}$ too.
\end{solution}

\begin{solution}{ex:crossing}
\idea Each crossing can be destroyed by deleting one edge.
\writeup Take a drawing with $\mathrm{cr}(G)$ crossings and delete one edge at each crossing; at least $m-\mathrm{cr}(G)$ edges remain, forming a plane graph, so $m-\mathrm{cr}(G)\le3n-6$.
For $K_6$, $m=15$, $n=6$: $\mathrm{cr}\ge3$. A drawing with $3$ crossings: put the vertices at $(0,0)$, $(6,0)$, $(3,5.2)$ (an outer triangle) and $(3.3,2.9)$, $(1.8,1.4)$, $(3.9,0.9)$ (a
rotated inner triangle) and draw all edges straight; exactly three pairs of edges cross (the notebook checks it).
\end{solution}

\begin{solution}{ex:hakimi}
\idea An orientation with out-degrees at most $3$ assigns each edge to one of its ends, each end used at most three times: a matching problem.
\writeup Form a bipartite graph with the edges of $G$ on one side and three copies $(v,1),(v,2),(v,3)$ of each vertex on the other, the edge $uv$ being adjacent to the six copies of $u$
and $v$. A matching covering all edges orients each edge away from the vertex it is matched to, with out-degrees at most $3$. Hall's condition: a set $F$ of edges with
$s$ ends has $3s$ neighbours; the graph formed by $F$ is simple and planar, so $|F|\le3s-6<3s$ if $s\ge3$, and $|F|=1<6$ if $s=2$.
\end{solution}

\begin{solution}{ex:greedyorder}
\idea Greedy never does worse than the colouring whose classes it follows.
\writeup Let $C_1,\dots,C_\chi$ be the classes of an optimal colouring and list $C_1$, then $C_2$, and so on. By induction, greedy gives each vertex of $C_j$ a colour at most $j$: its
earlier neighbours are in $C_1,\dots,C_{j-1}$ (none in $C_j$), with colours at most $j-1$, so some colour at most $j$ is free. For $K_{n,n}$ minus the matching $a_ib_i$, use the order
$a_1,b_1,a_2,b_2,\dots$: by induction $a_i$ and $b_i$ both get colour $i$, as they are adjacent to all earlier $a_j,b_j$ of the other side, which use colours $1,\dots,i-1$. Yet
the graph is bipartite.
\end{solution}

\begin{solution}{ex:interval}
\idea Earlier neighbours of an interval form a clique with it.
\writeup Order the intervals by left endpoint. When $I=[a,b]$ is coloured, every earlier interval $J$ meeting $I$ starts at or before $a$ and ends at or after $a$, so it contains $a$.
These intervals and $I$ all contain $a$: a clique. So at most $\omega-1$ colours are forbidden, and greedy uses at most $\omega$; as $\chi\ge\omega$, it is optimal. For register allocation
this means that the number of registers needed for straight-line code is the maximum number of variables alive at the same time.
\end{solution}

\begin{solution}{ex:mycielski}
\idea The shadow vertex $u_i$ can replace $v_i$ in any colouring, and $w$ forces a colour that the shadows avoid.
\writeup Triangles: the $u_i$ are pairwise non-adjacent, so a triangle contains at most one of them, and not $w$ together with a $v_j$ ($w$ has no $v_j$ neighbour). A triangle $u_iv_jv_l$ would give
$v_j,v_l$ adjacent neighbours of $v_i$, a triangle in $G$. A triangle among the $v_j$ is in $G$. Colouring: colour $v_i$ and $u_i$ alike with a $\chi(G)$-colouring and $w$ with a new colour:
$\chi(M(G))\le\chi(G)+1$. Conversely, suppose $M(G)$ is coloured with $k=\chi(G)$ colours, $w$ having colour $k$. No $u_i$ has colour $k$. Recolour each $v_i$ of colour $k$ with the colour of $u_i$. For an
edge $v_iv_j$: they cannot both have had colour $k$; if $v_i$ had it, its new colour, that of $u_i$, differs from that of $v_j$, which is adjacent to $u_i$. So $G$ has a colouring with $k-1$
colours, a contradiction. Starting from $K_2$, the graphs $M(K_2)=C_5$, $M(C_5)$ (the Gr\"otzsch graph, $11$ vertices) and so on are triangle-free with chromatic numbers
$2,3,4,\dots$.
\end{solution}

\begin{solution}{ex:classtwo}
\idea A colour class is a matching, which leaves a vertex uncovered when $n$ is odd.
\writeup With $n$ odd, each colour class has at most $\frac{n-1}2$ edges, while there are $\frac{\Delta n}2$ edges, so at least $\frac{\Delta n}{n-1}>\Delta$ colours are needed; by Vizing's theorem,
exactly $\Delta+1$. For $K_n$ with $n$ odd, $\chi'=n$; with $n$ even, $\chi'=n-1$ (the round-robin schedule of a tournament with an even number of teams, $n-1$ rounds).
\end{solution}

\begin{solution}{ex:chromcycle}
\idea Deleting an edge of a cycle gives a path, contracting it a shorter cycle.
\writeup Let $L_n$ be the path on $n$ vertices, a tree, so $P(L_n,k)=k(k-1)^{n-1}$. By Proposition~\ref{prop:delcon}, $P(C_n,k)=P(L_n,k)-P(C_{n-1},k)=k(k-1)^{n-1}-P(C_{n-1},k)$ for $n\ge4$, with $P(C_3,k)=k(k-1)(k-2)=(k-1)^3-(k-1)$. If the
formula holds for $n-1$, $P(C_n,k)=k(k-1)^{n-1}-(k-1)^{n-1}-(-1)^{n-1}(k-1)=(k-1)^n+(-1)^n(k-1)$. At $k=-1$: $(-1)^nP(C_n,-1)=2^n-2$, and indeed the acyclic orientations of a
cycle are all $2^n$ orientations but the $2$ directed ones.
\end{solution}

\begin{solution}{ex:eulerc}
\idea Run the proof of Euler's formula again: delete edges on cycles until a forest remains.
\writeup Induction on the number of edges lying on cycles. If an edge $e$ lies on a cycle, its two sides are in different faces, which merge in $G-e$ (as in the proof
of Theorem~\ref{thm:euler2}); $m$ and $f$ both drop by $1$ and $c$ does not change. If no edge lies on a cycle, $G$ is a forest, with $m=n-c$ edges and one face
(Lemma~\ref{lem:theta}): $n-m+f=c+1$.
\end{solution}

\begin{solution}{ex:listfive}
\idea Split along a chord, or peel off the vertex $v_k$ next to $v_1$, keeping two of its colours in reserve and forbidding them to its inner neighbours, which have colours to spare.
\writeup Induction on $n$; for $n=3$, $v_3$ has three colours, at most two forbidden. If $C$ has a chord $v_iv_j$, it splits $G$ into two such graphs $G_1\ni v_1v_2$ and $G_2$ sharing the edge $v_iv_j$;
colour $G_1$ by induction, then $G_2$ with $v_iv_j$ as the precoloured edge. Without chords, let $v_1,u_1,\dots,u_r,v_{k-1}$ be the neighbours of $v_k$ in order; since the inner faces at
$v_k$ are triangles, $v_1u_1\cdots u_rv_{k-1}$ is a path, the $u_j$ are inner vertices (no chords), and $G-v_k$ is again of the required type, with outer cycle $v_1\cdots v_{k-1}u_r\cdots u_1$.
Choose two colours $\gamma,\delta$ in the list of $v_k$ other than the colour of $v_1$, and remove them from the lists of the $u_j$, which keep at least $3$ colours. Colour $G-v_k$ by
induction. Then $v_k$ has neighbours $v_1$ (colour neither $\gamma$ nor $\delta$), the $u_j$ (neither $\gamma$ nor $\delta$) and $v_{k-1}$ (at most one of them): one of $\gamma,\delta$ is free for $v_k$. Every
planar graph is a subgraph of such a graph (add edges inside faces), with all lists of size $5$ and any two adjacent outer vertices precoloured from their lists.
\end{solution}

\section*{Sources and further reading}

\begin{itemize}
\item R.~Wilson, \emph{Four Colors Suffice: How the Map Problem Was Solved}, Princeton University Press, 2002 (revised 2014): the whole story,
with De Morgan's letter. K.~Appel and W.~Haken, \emph{Every planar map is four colorable}, Illinois Journal of Mathematics 21 (1977), 429--567;
N.~Robertson, D.~Sanders, P.~Seymour and R.~Thomas, \emph{The four-colour theorem}, Journal of Combinatorial Theory B 70 (1997), 2--44; G.~Gonthier,
\emph{Formal proof: the four-color theorem}, Notices of the AMS 55 (2008), 1382--1393.
\item K.~Kuratowski, \emph{Sur le probl\`eme des courbes gauches en topologie}, Fundamenta Mathematicae 15 (1930), 271--283; C.~Thomassen, \emph{Kuratowski's
theorem}, Journal of Graph Theory 5 (1981), 225--241; R.~Diestel, \emph{Graph Theory}, 5th ed., Springer, 2017, Chapter 4, whose presentation we
followed; B.~Mohar and C.~Thomassen, \emph{Graphs on Surfaces}, Johns Hopkins University Press, 2001, for all topological details.
\item P.~J.~Heawood, \emph{Map-colour theorem}, Quarterly Journal of Mathematics 24 (1890), 332--338; R.~L.~Brooks, \emph{On colouring the nodes of a
network}, Proceedings of the Cambridge Philosophical Society 37 (1941), 194--197; L.~Lov\'asz, \emph{Three short proofs in graph theory}, Journal of
Combinatorial Theory B 19 (1975), 269--271; V.~G.~Vizing, \emph{On an estimate of the chromatic class of a $p$-graph} (in Russian), Diskretnyi Analiz 3
(1964), 25--30.
\item H.~Whitney, \emph{A logical expansion in mathematics}, Bulletin of the AMS 38 (1932), 572--579; R.~P.~Stanley, \emph{Acyclic orientations of graphs},
Discrete Mathematics 5 (1973), 171--178; M.~Beck and R.~Sanyal, \emph{Combinatorial Reciprocity Theorems}, AMS, 2018.
\item C.~Thomassen, \emph{Every planar graph is 5-choosable}, Journal of Combinatorial Theory B 62 (1994), 180--181: a whole paper of one page and a half.
\end{itemize}
